\documentclass[10pt,a4paper]{amsart}
\usepackage[margin=19mm]{geometry}
\usepackage{amsmath,amssymb,mathtools}
\usepackage[scr=rsfs,cal=euler]{mathalfa}
\usepackage{xfrac}
\usepackage[unicode,hidelinks,bookmarks=false]{hyperref}
\newcommand{\ie}{i.e.\ }
\newcommand{\cf}{cf.\ }
\newcommand{\resp}{respectively\ }
\newcommand{\Subsection}{\subsection}
\providecommand{\nobreakdash}{\nobreak\hbox{-}}
\setlength{\emergencystretch}{2em}
\multlinegap0pt
\usepackage[shortlabels]{enumitem}
\usepackage{amssymb}
\usepackage{tikz}
\usepackage{xcolor}
\usepackage{mathtools}
\usepackage{tcolorbox}
\newtheorem{thm}{Theorem}
\title{Hardy spaces of discrete holomorphic functions \\on the upper half-lattice}
\author{Eugenio Dellepiane, Alessandro Monguzzi, Matteo Monti}
\date{Source: arXiv:2607.17726v1 (CC BY 4.0)}
\begin{document}
\maketitle

\noindent\emph{Source and licence.} Hardy spaces of discrete holomorphic functions \\on the upper half-lattice. Version arXiv:2607.17726v1 (CC BY 4.0), distributed by arXiv under the Creative Commons Attribution 4.0 International licence, http://creativecommons.org/licenses/by/4.0/, as recorded in the arXiv licence metadata for this exact version and shown on the abstract page https://arxiv.org/abs/2607.17726v1. Full licence notice: \texttt{licenses/arxiv-2607.17726-CC-BY-4.0.txt}.\par\smallskip

\noindent\textit{Editorial scope.} Counted unit: the article's thm T:approximationdiscrete (source0.tex lines 1430--1442). The article's complete original proof of it is delivered with it (source0.tex lines 1448--1602, one \texttt{proof} environment of 155 lines, which the article places away from the statement). No mathematics is written, repaired or re-derived: the statement and the proof are the article's own lines, in the article's own notation, and no retained line is substituted or re-flowed.  Retained uncounted as the source-local context the proof needs: lines 829--831; lines 1390--1392; lines 1402--1404.  Nothing was omitted from any retained span, so the package carries no omission record.  No retained line contains a citation, so the wrapper carries no bibliography and no bibliography entry.  No retained line contains a figure environment, an image-inclusion command or drawing code, so no image is delivered and the package declares no asset dependency.  Editorial formatting only, in addition: an AMS \texttt{amsart} preamble that restates the article's own declarations and macros used by the retained lines, namely the environments \texttt{thm} on separate counters, together with the standard packages those lines need.  The retained environments print in this document as Theorem 1 in order of appearance, whereas the article prints the same items with the numbering of its own full document.  The complete native source of the article as distributed in the arXiv e-print for this exact version is bundled unchanged as \texttt{sources/205592/source0.tex}.
    \begin{equation}\label{Eq:symmetryphi}
    \frac{\cos (\pi-t)}{1+\sin (\pi-t)}=-\frac{\cos t}{1+\sin t}, \qquad -\frac{\pi}{2}<t<\frac{3}{2}\pi,
    \end{equation}

    \begin{equation}\label{Eq:PWhalfplane}
    f(x+iy)=\frac{1}{\sqrt{2\pi}} \int_0^\infty g(t) e^{-yt} e^{ixt}\operatorname{d}\! t, \qquad x+iy\in\mathcal{U},
    \end{equation}

    \begin{equation}\label{eq:hvarepsilon}
    h(\varepsilon m+i\varepsilon n)=\frac{1}{\sqrt{2\pi}}\int_0^{\frac\pi\varepsilon}g(t)e_{\varepsilon t}(m+in)\operatorname{d}\!t= \frac{1}{\sqrt{2\pi}}\int_0^{\frac\pi\varepsilon} g(t) e^{im\varepsilon t}\Big(\frac{\cos(\varepsilon t)}{1+\sin(\varepsilon t)}\Big)^{n} \operatorname{d}\!t.
    \end{equation}

    \begin{thm}\label{T:approximationdiscrete}
    Let $f$ be as in \eqref{Eq:PWhalfplane}, $0<\varepsilon\leqslant1$ and $h$ as in \eqref{eq:hvarepsilon}. Then, for every compact set $K\subseteq \mathcal U$ we have that
    \[
    %\sup_{\substack{m,n:\\
    %\varepsilon m+i\varepsilon n\in K}}
    \sup_{K\cap\varepsilon\mathbb{Z}^2_+}|f-h|
    %\left|
    %f(\varepsilon m+i\varepsilon n)-h(\varepsilon %m+i\varepsilon n)
    %\right|
    \leqslant C_K \|g\|_{L^2((\frac{\pi}{2\varepsilon},\frac{\pi}{\varepsilon}))}+O(\varepsilon^{\frac{1}{4}}),
    \]
    where $C_K$ is a constant that only depends only on the compact $K$, not on $\varepsilon$ or $f$.
    \end{thm}

    \begin{proof}[Proof of Theorem \ref{T:approximationdiscrete}]
    There exist $0<a<b<\infty$ such that
    \(
    a\leqslant\operatorname{Im}z\leqslant b
    \) for every $z$ in $K$. Hence, whenever $\varepsilon m+i\varepsilon n\in K$, we  have 
    \begin{equation}\label{eq:bounds_n}
    0<a\leqslant\varepsilon n\leqslant b.
    \end{equation}
    We denote
    \[
    \varphi(t)=\frac{\cos(t)}{1+\sin(t)}, \qquad t\in[0,\pi].
    \]
    For $m,n$ such that $\varepsilon m+i\varepsilon n\in K$ we have that
    \begin{align*}
    \sqrt{2\pi}\big(f(\varepsilon m+i\varepsilon n)-h(\varepsilon m+i\varepsilon n)\big)
    &= \int_0^{\frac{\pi}{2\varepsilon}}
    g(t)
    (
    e^{-\varepsilon nt}-\varphi(\varepsilon t)^n
    )e^{i\varepsilon mt}\operatorname{d}\!t\\
    &\quad +  \int_{\frac{\pi}{2\varepsilon}}^{+\infty} g(t) e^{-\varepsilon n t }e^{i\varepsilon m t}\operatorname{d}\!t-
    \int_{\frac{\pi}{2\varepsilon}}^{\frac\pi\varepsilon} g(t)\varphi(\varepsilon t)^ne^{i\varepsilon m t}\operatorname{d}\!t\\
    &=I_{\varepsilon,m,n}
    +
    I\!I_{\varepsilon,m,n} 
    +
    I\!I\!I_{\varepsilon,m,n} 
    .
    \end{align*}
    We begin with $I\!I_{\varepsilon,m,n}$ and $
    I\!I\!I_{\varepsilon,m,n}$. Keeping in mind \eqref{eq:bounds_n}, we have
    \begin{align}\label{eq:II_bounds}
    \nonumber\sup_{\substack{m,n:\\
    \varepsilon m+i\varepsilon n\in K}} |I\!I_{\varepsilon,m,n}|&
    \leqslant \|g\|_{L^2((0,\infty))}\Big(\int_{\frac \pi{\varepsilon}}^{\infty}e^{-2\varepsilon n t}\operatorname{d}\!t\Big)^\frac{1}{2}\\
    &\leqslant \|g\|_{L^2((0,+\infty))}\Big(\int_{\frac \pi{\varepsilon}}^{+\infty}e^{-2a t}\operatorname{d}\!t\Big)^\frac{1}{2}=O(e^{-2a\frac{\pi}{\varepsilon}}),
    \end{align}
    for $0<\varepsilon<1$. About $
    I\!I\!I_{\varepsilon,m,n}$, by Cauchy--Schwarz,
    \begin{align}\label{Eq:IIIpt1}
        |I\!I\!I_{\varepsilon,m,n}|^2&\leqslant
        \|g\|_{L^2((\frac{\pi}{2\varepsilon},\frac{\pi}{\varepsilon}))}^2
        \int_{\frac{\pi}{2\varepsilon}}^{\frac{\pi}{\varepsilon}}
    |\varphi(\varepsilon t)|^{2n}
    \operatorname{d}\!t=
    \|g\|_{L^2((\frac{\pi}{2\varepsilon},\frac{\pi}{\varepsilon}))}^2\left(
      \frac1\varepsilon  \int_{\frac{\pi}{2}}^{\pi}
    |\varphi(t)|^{2n}
    \operatorname{d}\!t\right).
    \end{align}
    With the change of variable $s=\pi-t$, since $\phi$ enjoys the symmetry property in \eqref{Eq:symmetryphi}, we have that
    \begin{align}\label{Eq:IIIpt2}
      \frac1\varepsilon  \int_{\frac{\pi}{2}}^{\pi}
    |\varphi(t)|^{2n}
    \operatorname{d}\!t&= \frac1\varepsilon  \int_0^{\frac{\pi}{2}}
    |\varphi(s)|^{2n}
    \operatorname{d}\!s\leqslant\frac1\varepsilon  \int_0^{\infty}
    e^{-2ns}
    \operatorname{d}\!s = \frac{1}{2n\varepsilon}\leqslant\frac{1}{2a}.
    \end{align}
    We used the inequality $0\leqslant\phi(s)\leqslant e^{-s}$ for $0\leqslant s\leqslant\frac{\pi}{2}$. One can check it by differentiating.
    
    We turn to $I_{\varepsilon,m,n}$.  By Cauchy--Schwarz,
    \[
    \left|\int_0^{\frac{\pi}{2\varepsilon}}
    g(t)
    (
    e^{-\varepsilon nt}-\varphi(\varepsilon t)^n
    )e^{i\varepsilon mt}\operatorname{d}\!t\right|
    \leqslant
    \|g\|_{L^2((0,\infty))}
    \Big(
    \int_0^{\frac{\pi}{2\varepsilon}}
    |e^{-\varepsilon nt}
    -\varphi(\varepsilon t)^n
    |^2
    \operatorname{d}\!t
    \Big)^{1/2}.
    \]
    Since $\varepsilon<1,$ we split the integral in the right-hand side in two parts:
    \begin{align*}
        \int_0^{\frac{\pi}{2\varepsilon}}
    |e^{-\varepsilon nt}
    -\varphi(\varepsilon t)^n
    |^2
    \operatorname{d}\!t=&
    \int_0^{\frac{\pi}{2\sqrt\varepsilon}} |e^{-\varepsilon nt}-\varphi(\varepsilon t)^n|^2\operatorname{d}\!t +\int_{\frac{\pi}{2\sqrt\varepsilon}}^{\frac{\pi}{2\varepsilon}}
    |e^{-\varepsilon nt}
    -\varphi(\varepsilon t)^n
    |^2
    \operatorname{d}\!t.
    \end{align*}
    By a change of variable, we rewrite the first summand as
    \begin{align*}
    \int_0^{\frac{\pi}{2\sqrt\varepsilon}} |e^{-\varepsilon nt}-\varphi(\varepsilon t)^n|^2\operatorname{d}\!t&=\frac1\varepsilon\int_0^{\frac{\sqrt\varepsilon\pi}{2}} |e^{- nt}-\varphi(t)^n|^2\operatorname{d}\!t.
    \end{align*}
    Checking the Taylor expansions, one can show that $e^{-t}-\varphi(t)=t^3+o(t^3)$, as $t\to 0^+$. In particular, 
    \[
    |e^{-t}-\varphi(t)|\leqslant Ct^3, \qquad t\in(0,\sqrt{\varepsilon}\pi),
    \]
    for some constant $C>0$ that does not depend on $t$ nor $\varepsilon$. 
    Hence, from the identity
    \[
    A^n-B^n=(A-B)\sum_{k=0}^{n-1} A^{n-k-1}B^k,
    \]
    and the fact that $|e^{-t}|,|\varphi(t)|\leqslant 1$ for $t\in [0,\pi]$, we conclude that
    \begin{equation}\label{eq:I_main}
    \int_0^{\frac{\pi}{2\sqrt\varepsilon}} |e^{-\varepsilon nt}-\varphi(\varepsilon t)^n|^2\operatorname{d}\!t\leqslant \frac{n^2}{\varepsilon}\int_0^{\frac{\sqrt\varepsilon\pi}{2}}|e^{-t}-\varphi(t)|^2\operatorname{ d}\! t\leqslant C^2\frac{b^2}{\varepsilon^3}\int_0^{\frac{\sqrt\varepsilon\pi}{2}}t^6\operatorname{d}\!t =O(\sqrt{\varepsilon}).
    \end{equation}
    In the last inequality we also used the fact that $\varepsilon n\leqslant b$. We move on to the last summand.
    
    We use again that 
    \(
    e^{-t}\geqslant \varphi(t)\geqslant 0
    \)
    for $t$ in $[0,\frac\pi2]$.  
    In particular, whenever $t\in[0,\frac{\pi}{2\varepsilon}]$, % $[\frac{\pi}{4\sqrt\varepsilon},\frac{\pi}{2\varepsilon}]$ 
    we have the estimate %$0\leqslant \varphi(\varepsilon t)\leqslant e^{-\varepsilon t}$. Moreover,  $a\leqslant \varepsilon n$. Hence,  
    \[
    0\leqslant \varphi(\varepsilon t)^n
    \leqslant e^{-\varepsilon nt}
    \leqslant e^{-at}.
    \]
    Thus,
    \[
    |
    e^{-\varepsilon nt}
    -
    \varphi(\varepsilon t)^n
    |^2
    \leqslant
    (e^{-\varepsilon nt}
    +
    \varphi(\varepsilon t)^n)^2
    \leqslant
    4e^{-2at},\qquad t\in[0,\frac{\pi}{2\varepsilon}],
    \]
    and we conclude that
    \begin{equation}\label{eq:I_tail}
    \sup_{a\leqslant \varepsilon n\leqslant b}\int_{\frac{\pi}{2\sqrt\varepsilon}}^{\frac{\pi}{2\varepsilon}}
    |e^{-\varepsilon nt}
    -\varphi(\varepsilon t)^n
    |^2
    \operatorname{d}\!t
    \leqslant
    4\int_{\frac{\pi}{2\sqrt\varepsilon}}^{+\infty}e^{-2at}\operatorname{d}\!t =O(e^{-a\frac{\pi}{\sqrt\varepsilon}}).
    \end{equation}
    
    From \eqref{eq:II_bounds}, \eqref{Eq:IIIpt1}, \eqref{Eq:IIIpt2}, \eqref{eq:I_main} and \eqref{eq:I_tail} we conclude that
    \[
    \sup_{\substack{m,n:\\
    \varepsilon m+i\varepsilon n\in K}}|f(\varepsilon m+i\varepsilon n)-h(\varepsilon m+i\varepsilon n)|\leqslant \frac{1}{\sqrt{2a}}\|g\|_{L^2((\frac{\pi}{2\varepsilon},\frac{\pi}{\varepsilon}))}+O(\varepsilon^{\frac{1}{4}}),
    \]
    hence the proof.
    \end{proof}

\end{document}
